\chapter{Boolean Analysis — Gate Merge}
%%% generated by the blueprint pipeline from TCSlib.BooleanAnalysis.LMN.GateMerge
\section{Overview}

This module provides infrastructure for merging two gate arrays, indexed by
$\mathrm{Fin}\,m_1$ and $\mathrm{Fin}\,m_2$, into a single array indexed by
$\mathrm{Fin}\,(m_1 + m_2)$. It records the two projection equations for the merged
array, lemmas showing that circuit reindexing reroutes evaluation correctly through the
merged gates, and lemmas establishing that DNF width, variable injectivity, and
no-duplication properties are preserved under merging.

%%%%%%%%%%%%%%%%%%%%%%%%%%%%%%%%%%%%%%%%%%%%%%%%%%%%%%%%%%%%%%%%%%%%%%%%%%%%%%%
\section{Declarations}

\begin{definition}[Merge of two gate arrays]
\lean{LMN.mergeGates}
\label{LMN.mergeGates}
\leanok
Given gate arrays $g_1 : \mathrm{Fin}\,m_1 \to \alpha$ and
$g_2 : \mathrm{Fin}\,m_2 \to \alpha$, the merged array
$\mathrm{mergeGates}\,g_1\,g_2 : \mathrm{Fin}\,(m_1+m_2) \to \alpha$ takes its values
from $g_1$ on indices $0,\dots,m_1-1$ and from $g_2$ on indices
$m_1,\dots,m_1+m_2-1$.
\end{definition}

\begin{lemma}[Left component of a merged gate array]
\lean{LMN.mergeGates_castAdd}
\label{LMN.mergeGates_castAdd}
\leanok
\uses{LMN.mergeGates}
\difficulty{2}
Let $g_1 : \mathrm{Fin}\,m_1 \to \alpha$ and $g_2 : \mathrm{Fin}\,m_2 \to \alpha$ be two
gate arrays, and let $g_1 \ast g_2 : \mathrm{Fin}\,(m_1+m_2) \to \alpha$ denote their
merge, which agrees with $g_1$ on the indices $0,\dots,m_1-1$ and with $g_2$ on the
indices $m_1,\dots,m_1+m_2-1$. Then for every $i \in \mathrm{Fin}\,m_1$, evaluating the
merge at the index $i$ regarded as an element of $\mathrm{Fin}\,(m_1+m_2)$ (i.e. the
same value $i < m_1$) returns $g_1\,i$.
\end{lemma}

\begin{lemma}[Merged array on right-hand indices]
\lean{LMN.mergeGates_natAdd}
\label{LMN.mergeGates_natAdd}
\leanok
\uses{LMN.mergeGates}
\difficulty{2}
Let $g_1 : \mathrm{Fin}\,m_1 \to \alpha$ and $g_2 : \mathrm{Fin}\,m_2 \to \alpha$ be
gate arrays, and let $i \in \mathrm{Fin}\,m_2$. Then the merged array of $g_1$ and $g_2$
takes the value $g_2(i)$ at the index $m_1 + i$.
\end{lemma}

\begin{lemma}[Left re-indexing preserves circuit evaluation]
\lean{LMN.reidx_eval_mergeGates_left}
\label{LMN.reidx_eval_mergeGates_left}
\leanok
\uses{BoolCircuit.Circuit, BoolCircuit.Circuit.eval, BoolCircuit.Circuit.reidx, BoolCircuit.Circuit.reidx_eval, LMN.mergeGates}
\difficulty{2}
Let $c$ be a Boolean circuit on $m_1$ variables, and let $g_1 : \mathrm{Fin}\,m_1 \to
\bbf_2$ and $g_2 : \mathrm{Fin}\,m_2 \to \bbf_2$ be assignments. Write $\iota :
\mathrm{Fin}\,m_1 \to \mathrm{Fin}\,(m_1+m_2)$, $i \mapsto i$, for the order-preserving
inclusion of the first block, and let $g$ be the merged assignment on $m_1 + m_2$
variables that agrees with $g_1$ on the first $m_1$ indices and with $g_2$ on the last
$m_2$. Then the circuit obtained by replacing each literal index $i$ of $c$ by
$\iota(i)$ (keeping its sign), evaluated under $g$, has the same value as $c$ evaluated
under $g_1$.
\end{lemma}

\begin{lemma}[Right embedding evaluates a circuit through the merged assignment]
\lean{LMN.reidx_eval_mergeGates_right}
\label{LMN.reidx_eval_mergeGates_right}
\leanok
\uses{BoolCircuit.Circuit, BoolCircuit.Circuit.eval, BoolCircuit.Circuit.reidx, BoolCircuit.Circuit.reidx_eval, LMN.mergeGates}
\difficulty{2}
Let $m_1, m_2$ be natural numbers, let $c$ be a Boolean circuit on $m_2$ variables, and
let $g_1 : \bbf_2^{m_1}$ and $g_2 : \bbf_2^{m_2}$ be Boolean assignments. Form the
merged assignment on $m_1 + m_2$ variables that takes the values of $g_1$ on indices
$0,\dots,m_1-1$ and the values of $g_2$ on indices $m_1,\dots,m_1+m_2-1$, and re-index
$c$ by the right embedding $i \mapsto m_1 + i$ from $\mathrm{Fin}\,m_2$ into
$\mathrm{Fin}\,(m_1+m_2)$. Then \[ \bigl(\mathrm{reidx}(c,\; i \mapsto m_1 +
i)\bigr)(\text{merged assignment}) = c(g_2), \] where each side denotes circuit
evaluation; that is, evaluating the re-indexed circuit under the merged assignment gives
the same Boolean value as evaluating $c$ directly under $g_2$.
\end{lemma}

\begin{lemma}[Width bound on the left block of a merged array]
\lean{LMN.mergeGates_width_left}
\label{LMN.mergeGates_width_left}
\leanok
\uses{DNF, DNF.width, LMN.mergeGates}
\difficulty{1}
Let $g_1$ and $g_2$ be families of DNF formulas over $n$ variables, indexed by
$\mathrm{Fin}\,m_1$ and $\mathrm{Fin}\,m_2$ respectively, and let $\ell$ be a natural
number. Suppose every formula $g_1(k)$ has width at most $\ell$. Form the merged array
of length $m_1+m_2$, which reproduces $g_1$ on the first $m_1$ indices and $g_2$ on the
remaining ones. Then for each $i \in \mathrm{Fin}\,m_1$, the entry of the merged array
at the position of $i$ inside the first block again has width at most $\ell$.
\end{lemma}

\begin{lemma}[Width bound on the right block of a merged formula array]
\lean{LMN.mergeGates_width_right}
\label{LMN.mergeGates_width_right}
\leanok
\uses{DNF, DNF.width, LMN.mergeGates}
\difficulty{1}
Let $g_1$ and $g_2$ be families of DNF formulas in $n$ variables, indexed by
$\{0,\dots,m_1-1\}$ and $\{0,\dots,m_2-1\}$ respectively, and let $g$ be their merge:
the family on $\{0,\dots,m_1+m_2-1\}$ that agrees with $g_1$ on the first $m_1$ indices
and with $g_2$ on the remaining ones, so that $g(m_1+i)=g_2(i)$. Suppose $l$ is a
natural number for which every formula $g_2(k)$ has width at most $l$. Then for each
index $i$ of the second family, the merged formula $g(m_1+i)$ has width at most $l$.
\end{lemma}

\begin{lemma}[Width bound preserved under concatenating families of formulas]
\lean{LMN.mergeGates_width}
\label{LMN.mergeGates_width}
\leanok
\uses{DNF, DNF.width, LMN.mergeGates}
\difficulty{2}
Let $g_1$ and $g_2$ be two families of DNF formulas over $n$ variables, indexed
respectively by $m_1$ and $m_2$ positions, and let $l$ be a natural number. Suppose
every formula in $g_1$ and every formula in $g_2$ has width at most $l$. Then every
formula in the merged family of $m_1 + m_2$ formulas — those of $g_1$ in order, followed
by those of $g_2$ — also has width at most $l$.
\end{lemma}

\begin{lemma}[Variable-injectivity is preserved by merging gate arrays]
\lean{LMN.mergeGates_varInj}
\label{LMN.mergeGates_varInj}
\leanok
\uses{DNF, LMN.mergeGates}
\difficulty{2}
Let $g_1$ and $g_2$ be two arrays of DNF formulas over $n$ variables, indexed
respectively by $\{0,\dots,m_1-1\}$ and $\{0,\dots,m_2-1\}$, and let $g$ be their merge:
the array indexed by $\{0,\dots,m_1+m_2-1\}$ that agrees with $g_1$ on the first $m_1$
indices and with $g_2$, shifted down by $m_1$, on the remaining $m_2$. Suppose that in
every formula of $g_1$ and in every formula of $g_2$, any two literals lying in one and
the same term whose variables coincide are in fact the identical literal. Then the
merged array $g$ inherits this property: for every index $k$ and any two literals
occurring in a common term of the formula $g(k)$ that share the same variable, the two
literals are equal.
\end{lemma}

\begin{lemma}[Merging arrays of DNF formulas preserves duplicate-free terms]
\lean{LMN.mergeGates_nodup}
\label{LMN.mergeGates_nodup}
\leanok
\uses{DNF, LMN.mergeGates}
\difficulty{2}
Fix $n \in \bbn$, and let $g_1$ and $g_2$ be two arrays of DNF formulas over $n$
variables, of lengths $m_1$ and $m_2$ respectively; write $g_1 \sqcup g_2$ for their
merge, the array of length $m_1+m_2$ whose first $m_1$ entries are the formulas of
$g_1$, in order, and whose remaining $m_2$ entries are those of $g_2$, in order. Suppose
that every term occurring in any formula of $g_1$, and every term occurring in any
formula of $g_2$, is a list of literals with no repetition. Then the same holds for the
merged array: for every index $k$ and every term $t$ of the formula at position $k$ in
$g_1 \sqcup g_2$, the literals of $t$ are pairwise distinct.
\end{lemma}
